\documentclass[runningheads]{llncs}

\usepackage[T1]{fontenc}
\usepackage{graphicx,booktabs}
\usepackage{amssymb,mathtools,amsmath}
\usepackage{mathbbol}
\usepackage{listings}
\usepackage{lstcustom}
\usepackage{tikz}
\usetikzlibrary{automata, positioning, arrows}
\usepackage{macros}
\usepackage{comment}
\usepackage{todonotes}
\usepackage{jmllistings}
\usepackage[inline]{enumitem}
\usepackage{etoolbox}
\usepackage{mathpartir}
\usepackage{stmaryrd,textcomp}
\usepackage{mathabx}

\newtoggle{extended}
\toggletrue{extended} 
\togglefalse{extended} 
\lstdefinelanguage{stipula}{%
  basicstyle=\ttfamily,keywords={stipula,assets,fields,agreement,now,if,else,false,true}
}

\def\stipula#1{\mbox{\lstinline[mathescape=true,language=stipula]!#1!}}
\def\jml#1{\mbox{\lstinline[mathescape=true,language={[JML]Java}]|#1|}}

% Enumerate listings without prepending section numbers.
% There surely is a cleaner way to do this, but it seems to work for now.
\AtBeginDocument{\renewcommand*{\thelstlisting}{\arabic{lstlisting}}}

\renewcommand{\todo}[1]{}

\newif \ifconf
\conftrue

\begin{document}

% \title{\hspace*{-.35cm}\mbox{Verification of legal contracts via Java/JML and KeY}}

\title{Deductive Verification of Legal Contracts}

\titlerunning{Deductive Verification of Legal Contracts}

\author{Reiner H\"ahnle\inst{1}\orcidID{0000-0001-8000-7613} \and
  Cosimo Laneve\inst{2}\orcidID{0000-0002-0052-4061}}

\authorrunning{R.~H\"ahnle, C. Laneve}

\institute{
  Department of Computer Science, TU Darmstadt, Germany\\
  \email{reiner.haehnle@tu-darmstadt.de}
  \and
  DISI, University of Bologna, Italy\\
  \email{cosimo.laneve@unibo.it}
  }

\maketitle              

\begin{abstract}%
  We enable deductive verification for {\Stipula}, a domain-specific
  language for legal contracts, by systematically translating
  contracts into {\Java} programs annotated with {\tt Java Modeling
    Language} specifications, and subsequently proving them using
a deductive verification tool.
  % 
  A central challenge of the translation lies in representing
  event-driven and time-dependent behaviour within a static
  specification framework. To address this, we introduce a dispatch
  table that records schedulable events together with their triggering
  times.
  % 
  The technique is presented for acyclic {\Stipula} contracts. We also
  extend the proposed technique to cyclic contracts, outlining the
  additional challenges and the conditions under which the approach
  remains applicable.

  \keywords{Formal Verification \and Stipula Language \and
    Translation-based Verification \and Deductive Verification.}
\end{abstract}

\section{Introduction}
\label{sec:introduction}
%!TEX root = newmain.tex

{\Stipula} is a domain-specific programming language designed to model
legal contracts with formally enforceable properties, particularly
those governing asset transfers, obligations, and time-dependent
rights~\cite{CrafaLSV23}.  Contracts written in {\Stipula} combine
declarative clauses with executable semantics, yielding models whose
behaviour evolves dynamically over time. To facilitate the contract
lifecycle, a supporting toolchain has been developed that has been
integrated into a unified workbench, enabling legal practitioners and
developers in writing, testing, debugging, and managing legal
contracts~\cite{Laneve2026,stipulaprototype}.  While this
infrastructure provides practical support for contract development and
experimentation, it currently lacks advanced deductive verification
capabilities.

This paper complements the existing {\Stipula} workbench with a formal
verification layer capable of establishing both partial and total
correctness properties of legal contracts.  Our approach relies on a
systematic translation of {\Stipula} contracts into {\Java} programs
annotated with {\tt Java Modeling Language} ({\JML}) specifications,
thereby reducing the verification of legal contracts to the discharge
of proof obligations within a mature and well-established formal
methods framework.  Concretely, we use {\KeY}~\cite{KeYTutorial24} as
a general-purpose deductive verification system.

A central challenge in this setting is achieving sufficient precision
in the representation of execution behaviour. Legal contracts
typically exhibit \emph{dynamic} features -- notably the scheduling of
events at future times -- while {\JML} specifications are inherently
\emph{static}.  Modelling time and deferred events within a deductive
verification framework, therefore, requires an explicit and carefully
structured representation of the contract's temporal evolution. Unlike
our preliminary study~\cite{iFM25}, where time and events were treated
in an abstract manner, the present work fully incorporates timed
events into the translation. This is challenging because the
underlying semantic models of {\Stipula} contracts are infinite-state
due to the unbounded progression of time and the generation of future
events.

To represent the behaviour of {\Stipula} using static method contracts, we exploit the 
finite-state structure defined by the contract's underlying automaton, in line 
with {\Stipula}'s state-based programming style. Starting from this automaton, 
we systematically construct all sequences of clauses that are legally admissible 
according to the contract specification. This construction yields an explicit 
representation of the contract's behavioural space, abstracting away from concrete 
executions while preserving the permitted control-flow structures. A key component 
of this process is the dispatch table, which records the events that can be scheduled 
together with their trigger times. This representation plays a central role in 
verification: the soundness and completeness of the properties that can be established depend 
directly on how precisely the admissible execution paths are captured.

We first present the technique for acyclic {\Stipula} contracts. In
this setting, we formally define the translation, describe the
construction of the dispatch table and the generation of execution
paths, and show how partial-correctness properties of contracts can be
reduced to proof obligations discharged automatically by {\KeY}. For
acyclic contracts, the behavioural space is finite and can be
represented with high precision.

We subsequently extend our approach to cover cyclic contracts. The
translation devised for the acyclic setting can be generalised to
contracts featuring \emph{separate cycles}: cycles are not
nested, even if multiple paths between their entry/exit states
are permitted. Here, each cycle is mapped to a bounded {\Java} loop
iteration, whose termination is enforced through an explicit counter
and a predetermined upper bound.  This encoding makes the verification
task possible for a deductive verification tool, while
deliberately introducing a controlled over-approximation of the
contract's potentially unbounded behaviour. As customary in deductive
verification, the presence of loop iterations requires suitable
invariants. Accordingly, while the structural aspects of the
translation are fully automated, the user may be required to provide
contract-specific invariants to enable {\KeY} to complete the proof
obligations generated by cyclic constructs.

Overall, the main contributions of this work are as follows:
%
\begin{itemize}
\item \emph{A principled encoding of time-dependent contractual
    behaviour.}  We devise a novel and adequate data structure, the
  \emph{dispatch table}, for representing time constraints and event
  scheduling within a deductive verification setting.
\item \emph{A systematic and extensible translation framework.}  We
  define a formal, fully specified translation from legal contracts
  written in {\Stipula} into {\Java} programs annotated with {\JML}
  specifications. The translation is not limited to acyclic
  interaction patterns, but also covers separate (non-nested) cycles, 
  thereby addressing a wide range of
  realistic contractual behaviour.
\item \emph{A sound verification methodology based on mainstream
    deductive technology.}  We establish the soundness of the
  translation, ensuring that properties verified on the generated
  {\JML}/{\Java} code by a state-of-the-art deductive verification
  system such as {\KeY} correctly reflect the semantics of the
  original contract.
\end{itemize}

The paper is organised as follows. Section~\ref{sec:background}
contains the necessary background about {\Stipula} and
{\KeY}. Section~\ref{sec:translation} defines the translation function
for {\Stipula} contracts; the translation covers both acyclic
contracts and those with cycles.  Section~\ref{sec:dispatchtable}
analyses the management of timed events with the dispatch table and
the generation of execution paths; Section~\ref{sec:cyclic-behavior}
extends the analysis to cyclic behaviour.
Section~\ref{sec:relatedworks} briefly reviews related efforts in
formal verification of legal contracts and
Section~\ref{sec:conclusion} concludes the paper and outlines
directions for future work.




%%% Local Variables:
%%% mode: latex
%%% TeX-master: "newmain"
%%% End:

 
\section{Background}
\label{sec:background}
%!TEX root = newmain.tex

To set the stage for our translation-based verification method, we
briefly review the relevant background. Section \ref{subsect:stipula}
presents the {\Stipula} language, emphasising its constructs for
assets, states, and timed events.  Section \ref{subsect:key}
introduces the {\KeY} system, which we employ as the verification
backend.

\subsection{{\Stipula}}
\label{subsect:stipula}

A {\Stipula} contract consists of a set of parties, states, assets,
fields, and a set of functions and events. The declaration of a
contract is defined in Figure~\ref{tab:stipula-syntax}, where
\stipula{C} is the name of the contract, $\vect{\tt h}$ and
$\vect{\tt x}$ are the \emph{assets} and \emph{fields}, respectively,
$\vect{\tt A}$ are the \emph{parties}.
%
\begin{figure}[t]
  {\small
    \begin{tabular}{l@{\quad}l}
      $
      \begin{array}{l}
        \stipula{stipula C \{}
        \\
        \qquad \stipula{assets} \; \vect{\tt h}
        \\
        \qquad \stipula{fields} \; \vect{\tt x} 
        \\
        \qquad \stipula{agreement(} \vect{\tt A} \stipula{) \{  }
        \\
        \qquad \qquad \vect{{\tt A}_1} \; : \; \vect{{\tt x}_1}
        \\
        \qquad \qquad \cdots 
        % \quad {\tt //} \; \bigcup_{i \in 1..n} \vect{{\tt A}_i} \subseteq \vect{{\tt A}}   \; , \quad  \bigcup_{i \in 1..n} \vect{{\tt x}_i} \subseteq\vect{{\tt x}} \; , \quad	\bigcap_{i \in 1..n} \vect{{\tt x}_i} = \varnothing
        \\
        \qquad \qquad  \vect{{\tt A}_n} \; : \; \vect{{\tt x}_n}
        \\
        \qquad \stipula{\}} \tostate \Qwithat  
        \\
        \qquad     F
        \\
        \stipula{\}}
      \end{array}
      $
      &
        $
        \begin{array}{l@{\quad}l@{\quad}l}
          \mathit{Functions} & F \; ::= & \zero  \quad | \quad   \Qwithat \,\A : \, {\f}(\vect{\tt y})[\vect{\tt k}]\,(E)\{
                                          \,S\, ; \, W\,\} \,\tostate\, \Qwithat'\,  F
          \\
          \mathit{Prefixes} & P \; ::= & E \, \send \, {\x} \quad | \quad E \send \, \A \quad | \quad E \lolli {\h}, {\h}' 
\\
                             & & |\quad E \lolli {\h} ,  \A 
          \\
          \mathit{Statements} & S \; ::= &  \zero  \quad | \quad P \; S \quad | \quad \ifte{E}{S}{S} \, S
          \\
          \mathit{Events} &  W \; ::= & \zero  \quad | \quad \stipula{now + t} \,\event\, \Qwithat\, \{ \, S \,\} \,\tostate\, \Qwithat'\,  W 
          \\
          \mathit{Expressions }& E \; ::= &  v \quad | \quad \stipula{now} \quad | \quad X    
                                            \quad | \quad   E \, \stipula{op} \, E \quad | \quad  \stipula{nop} \, E
%\\
%\mathit{Values} & v ::= &  n \quad |\quad \stipula{false} \quad |\quad \stipula{true} \quad | \quad  s %\quad | \quad \Time \quad |\quad a 
          \\
          \\
          \\
          \\
\end{array}
        $
    \end{tabular}
    \vspace{-.6cm}
  }
  \caption{\label{tab:stipula-syntax} Syntax of {\Stipula}}
  \vspace{-.5cm}
\end{figure}
%
The \stipula{agreement} construct declares the parties that set the
initial value of the fields. For example, if the agreement has three
parties $\A_1, \A_2, \A_3$ and the contract has two fields
$\x_1, \x_2$, if it declares $\A_1: \x_1$ and $\A_2, \A_3 : \x_2$ then
$\x_1$ will be set by $\A_1$ and $\x_2$ will be set upon agreement on
the value between $\A_2$ and $\A_3$.
%
When the agreement is concluded, the initial state of the contract is
set and the parties may invoke a function in $F$.

A \emph{function} $\f$ declared as
$\Qwithat\,\A :\,{\f}(\vect{\tt y})[\vect{\tt k}]\,(E)\{\,S;\,W\,\}\,\tostate\,\Qwithat'$ 
can be invoked by a party $\A$ if the contract is in state $\Q$ and the guard $E$ is
\emph{true}.  The names $\vect{\tt y}$ and $\vect{\tt k}$ are the
formal parameters of {\f}; they are kept separate because
$\vect{\tt y}$ are \emph{field} values while $\vect{\tt k}$ are
\emph{asset} quantities.

Function bodies are \emph{statements} $S$ followed by \emph{events}
$W$.  The former include value transfers, asset movements, conditional
logic, and field assignments. {\Stipula} distinguishes between
different transfer operations: field updates and messages use the
symbol $\rightarrow$ with the usual semantics of assignment, while
asset transfers, called \emph{moves}, use the linear operator $\lolli$
to emphasise the conservation semantics. For instance, an expression
like ``$1 \lolli \stipula{wallet}, \stipula{Seller}$'' denotes
exclusive transfer of a unit in {\tt wallet} to the {\tt Seller} and,
\emph{at the same time}, the decrease by 1 of {\tt wallet}.  In
contrast ``\stipula{wallet} $\rightarrow$ \stipula{Seller}'' specifies
that the value currently stored in \stipula{wallet} is communicated to
\stipula{Seller}, while leaving the state of \stipula{wallet}
unchanged.  It is common practice in {\Stipula} to use the
abbreviations $\h\lolli\h'$ and $\h\lolli\A$, which stand for
$\h\lolli\h,\h'$ and $\h\lolli\h,\A$, respectively.
% These abbreviations improve readability and conciseness by omitting
% redundant information while preserving the underlying semantic
% structure of the corresponding move operations.

Events
$\stipula{now + t} \,\event\, \Qwithat\, \{ \, S \,\} \,\tostate\,
\Qwithat'$ define a statement
$S$ to be executed if, \emph{after} \stipula{t} time units from the
current execution, the contract is in the state $\Q$. The expression \stipula{t}
is either a non-negative number or a variable (a formal parameter or a 
field) or a variable plus a non-negative number.
If the event is
executed, the contract will transit to the state $\Q'$. The expression
\stipula{now + t} is called \emph{time guard}.  Functions and events
are generically called \emph{clauses}.

Expressions $E$ include integer constant values $v$, the keyword
\stipula{now} (that stands for a non-negative integer value), names of
assets, fields and parameters generically ranged over by $X$, standard
arithmetic and Boolean operations \stipula{op} (binary) and
\stipula{nop} (unary). Booleans are encoded by integers in the usual
way. Actually, {\Stipula} has a richer set of values (reals, strings, time and
asset values). For the time being, we stick to integer arithmetic,
where deductive verification systems exhibit a high degree of
automation. Other datatypes require corresponding support in automated
theorem provers and SMT solvers: experimentation is needed to
determine what is possible. We also want to keep the presentation simple.

% We refer to article~\cite{CrafaLSV23} for background on the design
% of {\Stipula}.
%
The semantics of {\Stipula} is defined operationally by means of a
transition relation between \emph{configurations}
$\contract {\comma} \Time$, where $\Time$ is a positive integer value
representing the system's global clock and $\contract$ is a tuple
$\C(\State\semi\ell\semi \Sigma \semi\SemEvent)$, where:
%
\begin{itemize}
\item $\C$ is the {\Stipula} contract name;
\item $\State$ is the contract's current state: either $\zero$ (for no
  state) or a contract state \state;
\item $\ell$ maps names (parties, fields, assets, function parameters)
  to values;
\item $\Sigma$ is the (possibly empty) residual of a function body or
  of an event handler, \emph{i.e.}~$\Sigma$ is either $\zero$ or a
  term $S\ W \tostate \texttt{@}\state$;
\item $\SemEvent$ is a (possibly empty) multiset of \emph{pending
    events} to be executed. We let $\SemEvent$ be $\zero$ when there
  are no pending events, or
  %
  \begin{equation}
    \label{eq:psi}
    {\Time}_1 \,\event_{\ev_1}\, \Qwithat_1\, \{ \, S_1 \,\}
    \,\tostate\, \Qwithat_1' \quad|~ \cdots ~|\quad {\Time}_n \,\event_{\ev_n}\,
    \Qwithat_n\, \{ \, S_n \,\} \,\tostate\, \Qwithat_n'\enspace,
  \end{equation}
  %
  where $\Time_i$ are the values of the expressions \stipula{now +}
  \stipula{t}$_i$ of the corresponding events.  The indices $\ev_i$
  are introduced by rule~\rulename{function} when the body of a
  function is incorporated into the configuration and denotes a unique
  identifier
  % . Each such index denotes the unique integer identifier
  statically assigned to the event.
\end{itemize}

The transition relation of {\Stipula} is
$\contract {\comma} \Time \lred{\mu} \contract' {\comma} \Time'$,
where $\mu$ may be empty, or an initial agreement, or
$A:\f(\wt{u})[\wt{v}]$ (in case of function invocations) or an event.
%
Its formal definition is given in Table~\ref{tab:controlrules} for the
control transition rules, see \cite{CrafaLSV23} for the agreement and
statement rules. The former rule set requires a specific
{\JML}/{\Java} encoding, whereas the latter are translated in a
standard way and are inessential for the understanding of this paper.
%
\begin{table}[t]
{\small
$
\begin{array}{c}
\mathrule{Function}{
	\begin{array}{c}
	\texttt{@}\state\,\A\,\texttt{:}\,\f\texttt{(} \wt{\mbox{{\tt y}}} \texttt{)[} \wt{\mbox{{\tt k}}} \texttt{]\,(} E {\tt )}\,{\tt \{}\,S\,W\,{\tt \}} \tostate \texttt{@}\state' \in \C
	\\
	\nored{\SemEvent}{\Time}
	%\quad \vect{v} \geq 0 
	\\
	\ell(\A) = A
	\quad
	\ell' = \ell [ \wt{\mbox{{\tt y}}} \mapsto \wt{u}, \wt{\mbox{{\tt k}}} \mapsto \wt{v} ]
	\quad
	\sem{E}_{\ell'} \neq 0
	\end{array}
	}{
	\C(\state \semi \ell \semi \zero \semi \SemEvent) \comma \Time 
	    \lred{A : \texttt{f}(\wt{u})[\wt{v}]}
	\C(\state \semi \ell' \semi S\,W \tostate \texttt{@}\state' 
	\semi \SemEvent) \comma \Time
}	
\quad
\mathrule{State-Change}{
    \sem{W\subst{\Time}{\texttt{now}}}_{\ell} = \SemEvent'
    }{
    \begin{array}{l}
	\C(\state \semi \ell \semi \zero \,W\tostate \texttt{@}\state' \semi \SemEvent) 
	\comma \Time 
	\\
	\qquad \lred{}
	\C(\state' \semi \ell \semi \zero \semi \SemEvent' ~|~\SemEvent) 
	\comma \Time
	\end{array}
}
\\
\\
\mathrule{Event-Match}{
	\begin{array}{c}
	\SemEvent = \Time \event_{\ev_\kappa} \texttt{@}\state~\{~S~\} \tostate \texttt{@}\state' ~|~ \SemEvent'
	\end{array}
	}{
	\begin{array}{l}
	\C(\state \semi \ell \semi \zero \semi \SemEvent) \comma \Time 
%	\\
%	\qquad 
\lred{\ev_\kappa}
	\C(\state \semi \ell \semi S \tostate \texttt{@}\state' \semi \SemEvent') 
	\comma \Time  
	\end{array} 
}
\quad
\mathrule{Tick}{
	\nored{\SemEvent}{\Time}
	}{
	\begin{array}{l}
	\C(\state \semi \ell \semi \zero \semi \SemEvent) \comma \Time 
	\\
	\qquad 
\lred{} 
	\C(\state \semi \ell \semi \zero \semi \SemEvent\downarrow_{\Time}) \comma \Time + 1
	\end{array}
	}
\\
\\
\end{array}
$
}
\caption{\label{tab:controlrules} The control transition rules of {\Stipula}}
\vspace{-.6cm}
\end{table}

The formal definition of
$\contract {\comma} \Time \lred{\mu} \contract' {\comma} \Time'$ uses
the following auxiliary predicates and functions:
%
\begin{itemize}
\item $\sem{E}_{\ell}$ is a \emph{partial function} returning the
  value of $E$ in memory $\ell$.
  % \item $\sem{E}_\ell^\mathsf{a}$ is the partial function returning
  %   asset values. This evaluation differs from $\sem{E}_{\ell}$
  %   because it never return negative values (assets are always
  %   nonnegative)
\item $\sem{W\{^\Time/_{{\tt now}}\}}_\ell$ is the multiset of
  scheduled events obtained from the sequence $W$ by replacing every
  {\tt now} by $\Time$ and evaluating every expression in the time
  guards.
\item $\nored{\SemEvent}{\Time}$ is a predicate that is \emph{true}
  whenever $\SemEvent$ is as in equation~\eqref{eq:psi}
  % $\SemEvent = \Time_1 \event \texttt{@}\state_1 \{ S_1 \} \tostate
  % \texttt{@}\state_1' ~|~ \cdots ~|~ \Time_n \event \texttt{@}\state_n
  % \{ S_n \} \tostate \texttt{@}\state_n'$
  and, for every $1 \leq i \leq n$, $\Time_i \neq \Time$. It is
  \emph{false} otherwise.
\item
$\SemEvent\downarrow_{\Time}$ removes from $\Psi$ every event with time value less or equal to $\Time$. 
\end{itemize}

We explain the rules in Table~\ref{tab:controlrules}.
%
Rule \rulename{Function} defines invocations: the label specifies
party $A$ performing the invocation of a function named {\tt f} with
the actual parameters. The transition may occur provided (\emph{i})
the contract is in the state \texttt{Q}, admitting invocations of {\tt
  f} from $A$, (\emph{ii}) the contract is \emph{idle},
\emph{i.e.}~the contract has no statement to execute (presence of
$\zero$ in the left configuration), (\emph{iii}) the precondition $E$
is satisfied, and no event can be triggered (premise
$\nored{\SemEvent}{\Time}$). In particular, the final constraint
expresses that events \emph{have precedence} on possible function
invocations. For example, if a payment deadline is reached and, at the
same time, the payment arrives, it will be refused in favour of the
event managing the deadline.  We use two notations for party
identifiers: {\A} and $A$. The former denotes the formal name used in
the contract specification, whereas $A$ refers to the concrete party
name instantiated and fixed when the agreement is concluded.

Rule \rulename{State-Change} says that a contract changes state when
the execution of the statement in the function body terminates. At
this stage, the sequence of events $W$ is added to the multiset of
pending events once their time expressions have been evaluated, in
particular \texttt{now} is replaced by the current value of the clock.
%
Rule \rulename{Event-Match} specifies that event handlers may run
provided there is no statement to perform and the time guard of the
event has exactly the value of $\Time$.
%
Rule \rulename{Tick} defines the elapsing of time, which may happen
when there is no statement to perform and no event can be triggered.

In this paper we use the state transition models of {\Stipula}
contracts, called the \emph{underlying automata}. The states of these
automata are those of the {\Stipula} contract. The transitions
correspond to clauses and are labelled either with the function name
(the party name is always omitted, for simplicity's sake we assume
that function names are pairwise different) or with a uniquely
numbered event. By definition, these automata are \emph{finite models};
% line number (\emph{e.g.}, \stipula{ev}$_{10}$ is the event at code
% line 10).
Figures~\ref{fig:license} and~\ref{fig:deposit-fsm} show the two
underlying automata of the contracts \stipula{Licence} and
\stipula{Deposit}.
%
These contracts are drawn from~\cite{iFM25} with minor modifications, which 
are discussed below. 

\begin{example}
  Figure~\ref{fig:license} defines the \stipula{License} contract
  regulating a licensing transaction between a \stipula{Licensor} and
  a \stipula{Licensee}, with a 30 days trial period and the
  possibility to purchase or decline the license. The field {\tt code}
  stores the information for the trial period access, the asset {\tt
    token} grants the licence.  If the \stipula{Licensee} does not
  purchase the license before the trial period expires, the contract
  terminates and the {\tt token} is automatically returned to the
  \stipula{Licensee}.
  % 
  The \stipula{agreement} in Line~\ref{line:lic:agb} allows the
  parties to agree on the license \stipula{cost}. Upon agreement, the
  first state is \stipula{@Init} where
  % 
\begin{figure}[h]
  % \centering
  \quad\quad
\begin{tabular}{@{}c@{\hspace{-.8cm}}c@{}}
{\scriptsize
\begin{lstlisting}[numbers=left,frame=none,xleftmargin=2.5em,mathescape=true,language=stipula,
escapechar=§]
stipula License {
    assets balance, token
    fields cost, code
    agreement (Licensor, Licensee)(cost) { §\label{line:lic:agb}§
        Licensor, Licensee : cost
    } $\tostate$ @Init §\label{line:lic:age}§
    @Init Licensor: offer(x)[n] {
        n $\lolli$ token
        x $\send$ code
        now + 1 $\event$ @Prop { token $\lolli$ Licensor } $\tostate$ @End §\label{line:event_init}§
    } $\tostate$ @Prop 
    @Prop Licensee: activate()[b] (b == cost) {
        b $\lolli$ balance
        code $\send$ Licensee ;
        now + 30 $\event$ @Trial {
                balance $\lolli$ Licensee
                token $\lolli$ Licensor
                } $\tostate$ @End
    } $\tostate$ @Trial
    @Trial Licensee: buy()[] {
        balance $\lolli$ Licensor
        token $\lolli$ Licensee
    } $\tostate$ @End
}
\end{lstlisting}}
&
\scalebox{0.80}{
\begin{tikzpicture}[->, >=stealth', node distance=2.5cm, every state/.style={draw, minimum size=1.2cm},initial text=]
  \node[state, initial] (init) {\stipula{Init}};
  \node[state, right of=init] (first) {\stipula{Prop}};
  \node[state, right of=first] (run) {\stipula{Trial}};
  \node[state, below of=first] (end) {\stipula{End}};

  \path (init) edge node[above] {\stipula{offer}} (first)
        (first) edge node[above] {\stipula{activate}} (run)
        (first) edge node[left] {\stipula{ev}$_{1}$} (end)
        (run) edge node[left] {\stipula{buy}} (end)
        (run.south) edge node[right] {\stipula{ev}$_{2}$} (end.east);
\end{tikzpicture}
}
\end{tabular}
%\vspace{-.3cm}
\caption{The License contract in {\Stipula} and its underlying automaton.}
\label{fig:license}
\vspace{-.3cm}
\end{figure}
%
the \stipula{Licensor} may invoke \stipula{offer}, transferring a
\stipula{token} (representing the license) into escrow and a trial
\stipula{code}.  A scheduled \stipula{event} is simultaneously
registered after 1 day to reclaim the token if the \stipula{Licensee}
fails to start the trial period.  The contract then moves to
\stipula{Prop}.
% 
If no function is called in this state, then time is advanced,
eventually triggering the event at line~\ref{line:event_init}.
% 
In \stipula{Prop}, the \stipula{Licensee} can activate the trial by
paying the \stipula{cost} (observe that \stipula{b} is an asset of the
\stipula{Licensee}), which is transferred to the contract's
\stipula{balance}: the contract acts as a notary for assets that are
not finally disposed.  Now the license code is revealed to the
\stipula{Licensee} and another \stipula{event} is scheduled to handle
the expiration of the trial period: if the license is not purchased
before 30 days, the balance is refunded, the token returned to the
\stipula{Licensor}.
%
The license purchase is done by the \stipula{buy} function, which
finalises the transaction by transferring the \stipula{balance} to the
\stipula{Licensor} and assigning the \stipula{token} permanently to
the \stipula{Licensee}. 
In the corresponding contract in~\cite{iFM25}, the start and end of the 
trial period are defined during the agreement, whereas here, for simplicity, 
they are hard-coded in the contract.
\end{example}

\begin{example}
  Figure~\ref{code:deposit-stipula} defines the \stipula{Deposit}
  contract that
%
\begin{figure}[t]
\centering
\begin{tabular}{c@{\hspace{-.6cm}}c}
{\scriptsize
\begin{lstlisting}[numbers=left,frame=none,xleftmargin=2.5em,mathescape=true,language=stipula,escapechar=§]
stipula Deposit {
    asset flour
    field cost_flour
    agreement (Client, Farm)(cost_flour) {
        Client, Farm : cost_flour
    } $\tostate$ @Start
    @Start Farm : begin()[h]{
        h $\lolli$ flour;
        now + 365 $\event$ @Run { flour $\lolli$ Farm } $\tostate$ @End
    } $\tostate$ @Run
    @Run Farm : send()[k]{
        k $\lolli$ flour
    } $\tostate$ @Run
    @Run Client : buy()[w](w/cost_flour <= flour){
        (w/cost_flour) $\lolli$ flour, Client
        w $\lolli$ Farm
    } $\tostate$ @Run
}
\end{lstlisting}}
& 
\scalebox{0.8}{
\begin{tikzpicture}[->, >=stealth', node distance=2.5cm, every state/.style={draw, minimum size=1.2cm}, initial text=]
  \node[state, initial] (start) {\stipula{Start}};
  \node[state, right of=start] (run) {\stipula{Run}};
  \node[state, right of=run] (end) {\stipula{End}};

  \path (start) edge node[above, sloped] {\stipula{begin}} (run)
        (run) edge[loop above] node[above] {\stipula{buy}} (run)
        (run) edge[loop below] node[below] {\stipula{send}} (run)
        (run) edge node[above, sloped] {\stipula{ev}$_{1}$} (end);
\end{tikzpicture}
}
\end{tabular}
%\vspace{-.3cm}
\caption{The \stipula{Deposit} contract in {\Stipula} and its underlying automaton}
\label{code:deposit-stipula}
\label{fig:deposit-fsm}
\vspace{-.5cm}
\end{figure}
%
models the interactions between a \stipula{Farm} and a
\stipula{Client}.  The \stipula{Farm} deposits flour, while the
\stipula{Client} purchases and withdraws the corresponding amount at
an agreed price.  The contractual terms are enforced over a validity
period of 365 days.  The clause \stipula{send} allows the
\stipula{Farm} to deposit flour into the contract's stock; the
operation
% \stipula{h $\send$ Client} is informs the \stipula{Client} about the
% deposited amount, while
\stipula{h $\lolli$ flour} implements the deposit in the asset
\stipula{flour}.  Actual delivery occurs only with \stipula{buy()[w]},
where \stipula{(w/cost_flour) $\lolli$ flour, Client} transfers flour
from the contract's stock to the \stipula{Client} (hence
\stipula{w/cost_flour} must not be greater than \stipula{flour}) and
\stipula{w $\lolli$ Farm} represents the payment to the
\stipula{Farm}.
% In this way, {\tt send} supplies the stock, whereas {\tt buy}
% withdraws from it under payment.
In~\cite{iFM25}, the \stipula{Deposit} contract does not allow \stipula{buy} and \stipula{send} 
to start and end in the same state, as the underlying theory does not cover this case. 
Consequently, a more convoluted contract was required. In contrast, the definition 
of \stipula{Deposit} presented here is simpler, as it relies on a more expressive 
theoretical framework.
\end{example}

\stipula{License} and \stipula{Deposit} represent paradigmatic
interaction patterns that frequently arise in legal contracts and that
are clearly reflected in their underlying automata. The behavioural
structure of \stipula{License} is \emph{acyclic}: no state is
revisited during execution.  In contrast, \stipula{Deposit} exhibits a
\emph{cyclic} pattern, as it includes functions that may be invoked
repeatedly, leading to repeated traversals of certain states.  It is
important to observe that, while the semantics of {\Stipula} admits
potentially infinite cyclic behaviour (for example, \stipula{Deposit}
allows an unbounded number of interactions in a fixed time window of
365 days), in practice only finitely many transactions can
realistically occur. In Section~\ref{sec:cyclic-behavior}, we use this
observation to obtain bounded encoding of cycles
% and simpler loop invariants, which are used to implement cyclic behaviour
in the translated contracts.

\begin{definition}[Cyclic Contract, Separate Cycle]
  \label{def:determinate}
  A {\Stipula} contract is \emph{cyclic} if its underlying automaton
  has a \emph{cycle}: a sequence
  $\Q_1 \lred{\nu_1} \Q_{2} \lred{\nu_2} \cdots \lred{\nu_{\ell-1}}
  \Q_\ell \lred{\nu_\ell} \Q_{1}$ where $\Q_1,\ldots,\Q_\ell$ are
  pairwise different and $\nu_i$ are either function names or
  events. A contract has \emph{separate cycles} if any two cycles
  either have no state in common or share their initial state only.
  % $\Q_1 \lred{\mu_1} \Q_{2} \lred{\mu_2} \cdots \lred{\mu_{\ell-1}} \Q_\ell 
  % \lred{\mu_\ell} \Q_{1}$
  % of pairwise different states $\Q_1,\ldots,\Q_n$ with transitions
  % $\Q_i \lred{\mu_i} \Q_{i+1}$ and $\Q_n \lred{\mu_n} \Q_{1}$ where
  % $\mu_i$ are either function names or events.  The contract has
  % \emph{disjoint cycles} if different cycles have no state in common.
  A contract is \emph{acyclic} if it has no cycle.
\end{definition}

The contract \stipula{License} is acyclic, while \stipula{Deposit} has
separate cycles.  In what follows, we first consider acyclic contracts
and then turn to separate cyclic ones. This progression allows us to
introduce the complexities of the translation into {\JML}/{\Java} in a gradual
manner.

\subsection{Deductive Verification with the {\KeY} System}
\label{subsect:key}

The {\KeY} system \cite{DBLP:series/lncs/10001,KeYTutorial24} is a
deductive verification framework for {\Java} that combines symbolic
execution, invariant reasoning, and method contracts\footnote{Not to
  be confused with {\Stipula} contracts. In this paper, both kinds of
  contract are featured, but it should be obvious from the context
  when we mean {\Stipula} contracts and when {\JML}/{\Java} method
  contracts.} on top of a calculus for an expressive program logic.
%
The main use case of {\KeY} is formal verification of Java programs
annotated with specifications written in {\JML}. 
It provides an interactive user environment, where one can
construct correctness proofs of {\Java} methods against their {\JML}
contracts. Method contracts typically include preconditions
(\jml{requires}), postconditions (\jml{ensures}), frame
conditions (\jml{assignable}), class invariants, and auxiliary
annotations such as loop invariants and assertions. Given a
{\JML}-annotated {\Java} class, {\KeY} translates the specifications into
logical proof obligations and attempts to discharge them using its
symbolic execution engine.
%
The details of the verification process are irrelevant for the purpose
of this paper: {\KeY} is used as a \emph{black box}, the input is a
{\JML}-annotated Java file that results from translation of a
{\Stipula} contract. In principle, \emph{any} deductive verification
system with a sufficient coverage of \JML/\Java\ can be used, for
example, Krakatoa~\cite{Krakatoa07} or OpenJML~\cite{Cok21}, although
we did not try that.  In general, {\KeY} is used in interactive or
auto-active mode, however, for all case studies discussed below, the
verification is \emph{fully automatic}.

%%% Local Variables:
%%% mode: latex
%%% TeX-master: "newmain"
%%% End:


\section{Translation Approach}
\label{sec:translation}
%!TEX root = newmain.tex

To enable formal verification of a {\Stipula} contract with {\KeY}, we
define a systematic translation into {\Java} code annotated with
{\JML} specifications. As discussed in~\cite{iFM25}, {\Stipula}
contracts are translated into classes where functions and events are
modelled by stateful {\Java} methods, assets and fields are modelled
by suitable {\Java} fields.  The definition of the translation is
given in Listing~\ref{tab:translate}.

We begin by analysing the declaration component of the translation.
As mentioned in Section~\ref{subsect:stipula}, we assume all {\Stipula}
fields, assets, and auxiliary variables are of type \jml{int} which
preserves the operational semantics of the original
contracts.\footnote{{\Stipula} itself is type-free: concrete types
  (e.g., integers, doubles, strings, Booleans) are inferred statically
  and at run time~\cite{CrafaLSV23}.}
%
Verification involves a trade-off between precision and generality:
parameters may be assigned concrete values, thereby analysing a
specific execution scenario, or left symbolic to reason about all
possible executions.  The translation in Listing~\ref{tab:translate}
adopts the former approach, targeting a single scenario.
%
Accordingly, fields $\vect{{\tt x}}$ are initialised to constant
values $\vect{u}$ (the values the fields have after the agreement),
and formal parameters of functions are likewise instantiated to fixed
values
($\vect{\tt y_1} = \vect{v_1}, \vect{\tt k_1} = \vect{v_1'}, \ldots,
\vect{\tt y_n} = \vect{v_n}, \vect{\tt k_n} = \vect{v_n'}$).

{\small
\begin{lstlisting}[frame=none,numbers=none,language={[JML]Java},mathescape=true,caption={The {\sc translate} function for a {\Stipula} contract},label={tab:translate}]
$\mbox{\sc translate}$( stipula C { assets $\vect{\tt h}$  fields $\vect{\tt x}$   agreement( $\vect{\tt A}$ ) {  $\cdots$ } $\tostate$ $\Qwithat$   F  ) $\eqdef$
  public final class C {
     public static int $\vect{\tt A}$, $\vect{\tt h}$, $(A{\unds}h)^{A \in \vect{\tt A}, h \in \vect{\tt h}}$ ;
     public static int $\vect{\tt x}$ = $\vect{u}$ ;
                    // $\textcolor{commentColor}{\vect{u}}$ are the values the fields $\textcolor{commentColor}{\vect{x}}$ have after the agreement
     public static int now = 0 ;
     public static int u_Q1, ..., u_Qm, w_Q1, ..., w_Qm ; 
                    // auxiliary variables: Q1, ..., Qm are the states of C
     public static int[] DT = new int[N] ;
                    // DT is the dispatch table, N is the number of events in C   
     public static int $\vect{\tt y_1} = \vect{v_1}, \vect{\tt k_1} = \vect{v_1'}$,..., $\vect{\tt y_n} = \vect{v_n}, \vect{\tt k_n} = \vect{v_n'}$ ; 
                    // formal parameters of functions and their initialization                                       
     public static invariant $\assetInv{\C}$ ;
     public static invariant 1 <= w_Q1 <= max_Q1 && $\cdots$ && 1 <= w_Qm <= max_Qm ;
        // max_Q1,...,max_Qm are the number of functions starting at Q1,...,Qm
     public static invariant 0 <= u_Q1 <= max_tv && $\cdots$ && 0 <= u_Qm <= max_tv ;
                    // max_tv is a constant, see Section 4
     $\translate{F}$
     
     /*@ public normal_behaviour
       @ requires $\vect{\tt h}$ == 0 ; DT[0] == -1 && ... && DT[N-1] = -1 ;   
         // assets of C are initially 0, DT is empty (encoded by value -1)
       @ requires OTHER PRE-CONDITIONS ;
       @ ensures THE POST-CONDITION ;       // the property we want to prove
       @*/

     public final static void behaviour(){  $\genscenario{\C}{\Q}$  }
}
\end{lstlisting}
}

The translation introduces auxiliary variables: The dispatch table
\jml{DT} is an array of size \jml{N}, where \jml{N} is the number of
events in the contract.  For each contract state \jml{Q1}, $\ldots$,
\jml{Qm}, the variable \jml{u_Qi} tracks time progression and ranges
over non-negative values bounded by \jml{max_tv}, the largest constant
(plus one) appearing in time expressions (in principle, a single
variable would suffice).  For each state \jml{Qi}, the variable
\jml{w_Qi} counts function invocations and ranges between 0 and
\jml{max_Qi} (the number of functions whose initial state is
\jml{Qi}).

Assets satisfy a non-duplication property: \emph{the total amount of a
  given asset in the system} -- that is, in the contract and in all
the parties -- \emph{must remain constant}.  To enforce this property,
we introduce a field $A\_h$ for every party $A$ and every asset
\jml{h} in class $\C$. Moreover, we enforce a static invariant
$\assetInv{\C}$ that holds for every method in $\C$ as follows:

{\small
\begin{lstlisting}[frame=none,numbers=none,language={[JML]Java}, mathescape=true]
 $\assetInv{\C}$ $\eqdef$ 
      $h_1$ + $A_1{\unds}h_1$ + $\cdots$ + $A_k{\unds}h_1$ == \old($h_1$) + \old($A_1{\unds}h_1$) + $\cdots$ + \old($A_k{\unds}h_1$)
      && $\cdots$ &&
      $h_r$ + $A_1{\unds}h_r$ + $\cdots$ + $A_k{\unds}h_r$ == \old($h_r$) + \old($A_1{\unds}h_r$) + $\cdots$ + \old($A_k{\unds}h_r$)
\end{lstlisting}
} 

In \JML\ the \jml{old} construct is used to refer to the values of
fields \emph{prior to method execution}. This allows us to precisely
relate the post-state of a method to its pre-state, which is essential
for enforcing the non-duplication property.
%
For example, $\assetInv{{\sf License}}$, the invariant of the {\sf
  License} contract is

{\footnotesize
\begin{lstlisting}[frame=none,numbers=none,language={[JML]Java},mathescape=true]
   balance + Licensor_balance + Licensee_balance ==
       \old(balance) + \old(Licensor_balance) + \old(Licensee_balance)
&& token + Licensor_token + Licensee_token ==
       \old(token) + \old(Licensor_token) + \old(Licensee_token)
\end{lstlisting}
}

When assets are transferred directly from one party to another without
being stored in a contract's asset (for example, the \stipula{buy}
function in the \stipula{Deposit} contract), we also add the invariant

{\small
\begin{lstlisting}[frame=none,numbers=none,language={[JML]Java}, mathescape=true]
   $\A_1$ + $\cdots$ + $\A_k$ = \old($\A_1$) + $\cdots$ + \old($\A_k$)
\end{lstlisting}
} 

\noindent
modeling that, in these cases, assets are transferred directly to the
party's identifier (which was introduced for this
purpose).\footnote{In~\cite{iFM25}, assets are classified as
  \emph{divisible} (\emph{e.g.}, digital currencies) or
  \emph{indivisible} (\emph{e.g.}, non-fungible tokens), thus
  generating appropriate invariants to enforce the intended asset
  semantics. Here, for simplicity, we assume that every asset is
  divisible.}

Functions and events are modeled as static {\Java} methods annotated
with {\JML} specifications consisting of three clauses. Consider a
{\Stipula} function declared as
$ \Qwithat \,\A\,\texttt{:}\,\f\texttt{(} \wt{\tt y} \texttt{)[}
\wt{\tt k} \texttt{]\,(} E {\tt )}\,{\tt \{}\,S\,W\,{\tt \}} \tostate
\Qwithat' $. Its translation, defined in
Listing~\ref{tab:translatefun}, includes:
%
\begin{itemize}
\item A \jml{requires} clause expressing that guard $E$ that must hold
  in the initial state for the function to execute.
\item An \jml{ensures} clause specifying the resulting state
  transition, including field updates and asset transfers. This clause
  is defined via the strongest postcondition $\StrongPostCond{E,S}$
  computed from the pre-state (see~\cite{KeYTutorial24} for details
  about strong post-conditions).
\item An \jml{assignable} clause listing the variables that may be
  modified, determined by a simple static ``writes-to'' analysis
  $\LhsVar{\A}{S}$.
  % 
\end{itemize}

The body of each method corresponding to a {\Stipula} function invoked
by a party {\A}, as well as the bodies of events declared within that
function, are generated via the function $\StipulaToJava{\A}{S}$. This
function maps a {\Stipula} statement $S$ to a sequence of {\Java}
statements, yielding a compositional encoding of the function's
operational semantics into executable {\Java} code.

{\small
\begin{lstlisting}[frame=none,numbers=none,language={[JML]Java},mathescape=true,caption={The {\sc translate} function for {\Stipula} functions and events},label={tab:translatefun}]
$\mbox{\sc translate}$( $\Qwithat \,\A\,\texttt{:}\,\f\texttt{(} \wt{\tt y} \texttt{)[} \wt{\tt k} \texttt{]\,(} E {\tt )}\,{\tt \{}\,S\,W\,{\tt \}} \tostate \Qwithat'$ ) $\eqdef$
 /*@ public normal_behaviour
   @ requires $E$ ;
   @ ensures  $\StrongPostCond{E,S}$ && $\ensureDT_\f$ ; 
   @ assignable $\LhsVar{\A}{S}$, DT ;
   @*/
 public final static void f($\wt{{\tt int \; y}}$, $\wt{{\tt int \; k}} $) {
     $\StipulaToJava{\A}{S}$ 
     $\updateDT{}{W}$    
}
 /*@ public normal_behaviour        // for every $\textcolor{commentColor}{{\tt now}+{\tt t}_i \,\event\, \Qwithat_i\, \{ \, S_i \,\} \,\tostate\, \Qwithat_i' \in W}$
   @ ensures  $\StrongPostCond{{\tt true},S_i}$ ;
   @ assignable $\LhsVar{\A}{S_i}$ ;
   @*/
 public final static void event_i() {
     $\StipulaToJava{\A}{S_i}$ 
}
\end{lstlisting}
}

  The definition of $\StipulaToJava{\A}{S}$ follows a well-established
  and standard formulation. For readability, the definition is
  relegated to the Addendum~\ref{sec:stipula2java} at the end of the paper.
%

The discussion about the
% preamble clause $\ensureDT_\f$ and the
function $\genscenario{\C}{\Q}$ (in
Listing~\ref{tab:translate}) and the functions $\updateDT{}{W}$ and
\ensureDT$_\f$ (in Listing~\ref{tab:translatefun})
are deferred to the next section, as they
rely on the dispatch table mechanism.

%%% Local Variables:
%%% mode: latex
%%% TeX-master: "newmain"
%%% End:


\section{The Dispatch Table and Generation of Execution Paths}
\label{sec:dispatchtable}
%!TEX root = newmain.tex

To complete the translation, we encode the behaviour of a {\Stipula}
contract.  We use three key mechanisms:
%
\begin{enumerate}
\item Event declarations are collected into a \emph{dispatch table},
  recording enabled events and their execution times.
\item \emph{Symbolic execution} implicitly enumerates admissible
  scenarios via non-deter\-ministic function and event selection;
\item Time progression in a state {\Q} is bounded by a constant
  \stipula{max_tv}, defined as the maximum constant (plus one)
  occurring in time expressions, and tracked through the symbolic
  variable ${\tt u{\unds}Q} \in [0 .. \stipula{max_tv}]$. Advancing
  time beyond this bound is unnecessary, as no pending event can be
  executed past that deadline unless generated by a subsequent
  function invocation.
\end{enumerate}

  The dispatch table, introduced in Listing~\ref{tab:translate}, is an
  array \jml{DT} of size \jml{N}, where \jml{N} is the number of
  events in the contract.  The $i$-th entry corresponds to the $i$-th
  event
  $\stipula{now + t}_i\,\event\,\Qwithat_i\, \{ S_i \} \tostate\,
  \Qwithat_i'$.  The value \jml{DT[i]} is \jml{-1} if the event has
  not been created; otherwise, it stores the scheduled execution time
  $\jml{now + t}_i$, where \jml{now} denotes the time at which the
  event was generated.

  Events declared in a function are inserted in
  Listing~\ref{tab:translatefun} into the dispatch table by
  $\updateDT{\Q}{W}$ . If a function $\f$ declares events indexed by
  $\ev_{\jmath_1}, \ldots, \ev_{\jmath_{N_\f}}$ with time expressions
  $\jml{now + t}_{\jmath_i}$, then $\updateDT{\Q}{W}$ assigns the
  corresponding computed times to the entries
  \jml{DT[$\ev_{\jmath_i}$]}.  In other words, it updates exactly
  those positions associated with the events generated by
  $\f$. Entries of \jml{DT} not corresponding to events in $W$ remain
  unchanged, while the relevant entries are updated to
  $\jml{now + t}_{\jmath_i}$ using the current value of \jml{now}.
  The formalisation of this behaviour in {\JML} is done by the
  predicate \ensureDT$_\f$. Thus, the specification guarantees that
  precisely the events declared in $\f$ are scheduled, relative to the
  current time \jml{now}, and that no other dispatch table entries are
  affected.  (The formal definition of $\updateDT{\Q}{W}$ and
  {\ensureDT}$_\f$ is omitted, because it is straightforward.)


To obtain the possible behaviour of a \Stipula\ contract, the dispatch
table will be queried to identify one of the events that can be
enabled next. To this end, we introduce the auxiliary \Java\ method
%
\begin{center}
  \jml{int minTimeEntry(\{ev$_1$, $\ldots$ , ev$_m$\}) }\enspace.
\end{center}

% where {\tt ev}$_1$, $\cdots$ , {\tt ev}$_m$ are natural numbers
% encoding events of the contract.
The method returns the index in the dispatch table entry corresponding
to an event in \{{\tt ev}$_1$, $\ldots$ , {\tt ev}$_m$\} whose
associated time value is minimal, provided that such an event exists;
otherwise, it returns {\tt -1}. The method is non-deterministic, as
multiple events may share the same minimal time value in the dispatch
table. In this case, any one of these events may be selected, thereby
allowing the symbolic execution to explore all admissible scheduling
choices.  The implementation of the method \jml{minTimeEntry()} is
omitted. Instead, we provide its {\JML} specification that defines the
essential semantic properties required for verification, which is
sufficient for {\KeY} to reason symbolically about the method's
behaviour:

{\small
\begin{lstlisting}[language={[JML]Java},frame=none,numbers=none,mathescape=true,escapechar=']
/*@ public normal_behaviour // specification of int minTimeEntry({ev$_1$,$\ldots$, ev$_m$})
  @ ensures ( \exists i; 0 <= i < N; (i $\in$ {ev$_1$, $\ldots$, ev$_m$} && DT[i] >= now) )
            $\Rightarrow$ \result $\in$ {ev$_1$, $\ldots$, ev$_m$} && DT[\result] >= now && 
               \forall j; 0 <= j < N; (j $\in$ {ev$_1$, $\ldots$, ev$_m$} && DT[j] >= now) 
                                            $\Rightarrow$  DT[j] >= DT[\result] ;
  @ ensures ( \forall i; 0 <= i < N; i $\in$ {ev$_1$, $\ldots$, ev$_m$} $\Rightarrow$ DT[i] < now)
            $\Rightarrow$ \result == -1;
  @ assignable \nothing; 
  @*/
\end{lstlisting}
}

The first \jml{ensures} clause says that when there is an index in
\{\ev$_1$, $\ldots$, \ev$_m$\} such that its time value in \jml{DT} is
greater or equal to \jml{now}, then such an index is returned. There
can be more than one such index. This causes symbolic execution to
branch over all indices \jml{$\mathtt{\backslash}$result} such that
\jml{DT[$\mathtt{\backslash}$result] >= now}. The second \jml{ensures}
clause covers the case when no entry is found. In this case {\tt
  {\textbackslash}result} is \jml{-1}.

{\small
\begin{lstlisting}[frame=none,numbers=none,language={[JML]Java},mathescape=true,caption={The function $\genscenario{\C}{\Q}$; 
the ${\f}_1, \ldots, {\f}_n$ and ${\tt event\_1}, \ldots, {\tt event\_m}$ are the functions and events that start in {\Q} and end in $\Q_1, \ldots, \Q_n$ and 
$\Q_1', \ldots , \Q_m'$, respectively.},label={tab:fungenerate},float]
$\genscenario{\C}{\Q} \; \eqdef$ 
    int entry = minTimeEntry({ev$_1$,$\ldots$, ev$_m$}) ; 
    if ((entry == -1) || (DT[entry] > now && w_Q > 0)) { // call a function 
       if (entry == -1) now = now + u_Q ; else now = min(DT[entry]-1, now+u_Q) ;
       if (w_Q == 1) { $\f_1(\vect{\tt y_1}, \vect{\tt k_1})$ ; $\genscenario{\C}{\Q_1}$ }
       else if (w_Q == 2) { $\f_2(\vect{\tt y_2}, \vect{\tt k_2})$ ; $\genscenario{\C}{\Q_2}$ }
       ...
       else { $\f_n(\vect{\tt y_n}, \vect{\tt k_n})$ ; $\genscenario{\C}{\Q_n}$ } // max_Q == n, see Listing 1
    } else { // trigger an event
       now = DT[entry] ;
       if (entry == 1) { event_1() ; DT[entry] = -1 ; $\genscenario{\C}{\Q_1'}$ }
       else if (entry == 2) { event_2() ; DT[entry] = -1 ; $\genscenario{\C}{\Q_2'}$ }
       ...
       else { event_m() ; DT[entry] = -1 ; $\genscenario{\C}{\Q_m'}$ } // entry == m
    } 
\end{lstlisting}
}

  We can now complete the definition of \textsc{translate}
  (Listing~\ref{tab:translatefun}) by specifying the function
  $\genscenario{\C}{\Q}$ (Listing~\ref{tab:fungenerate}). Assume that
  in state $\Q$ functions ${\f}_1, \ldots, {\f}_n$ can be called and
  events ${\tt event}_1, \ldots, {\tt event}_m$ can be triggered. The
  {\Stipula} semantics is non-deterministic. To model this, we use two
  symbolic variables that were declared in
  Listing~\ref{tab:translate}: \jml{w_Q} $\in [\texttt{0}..\texttt{max\_Q}]$, for selecting a
  function in the state {\Q}, and \jml{u_Q} $\in [0..{\tt max\_tv}]$, for controlling time
  progression.  The behaviour is as follows:

  \begin{enumerate}
  \item \emph{Enabled events at current time}.  If
    \jml{minTimeEntry(...)} returns an index whose time equals
    \jml{now}, the corresponding event \emph{must} execute (events
    preempt functions). After execution, the event is removed from
    \jml{DT}. (This removal is actually unnecessary since the
    automaton is acyclic, it cannot occur again.)
  \item \emph{Pending future events.}  If one of the earliest
    scheduled events occurs at a time \jml{entry > now},
    then two cases arise: either time advances, but not until \jml{entry},
    i.e.\ at most to \jml{min(DT[entry]-1, now+u_Q)} and a function is
    executed (\jml{w_Q > 0}); otherwise, time advances \emph{exactly}
    to \jml{t} and the event is executed (\jml{w_Q = 0}).
  \item \emph{No pending events}.  If no event from $\Q$ remains to be
    scheduled, a function must eventually execute.  Time progression
    is bound by \jml{max_tv}, defined as the largest constant (plus
    one) appearing in time expressions. Assigning \jml{u_Q $\,\in\,$
      [0..max_tv]} ensures a finite, yet sound exploration of all
    relevant execution scenarios.
\end{enumerate}

Now that the translation is complete, it becomes possible to define and
verify meaningful properties for {\Stipula} contracts translated to
{\JML}/{\Java} by filling in the slots for \jml{OTHER PRE-CONDITION} and
\jml{THE POST-CONDITION}, respectively, in
Listing~\ref{tab:translate}. For example, one can prove with {\KeY}
the following property of the {\tt License} contract:

{\footnotesize
\begin{lstlisting}[language={[JML]Java},frame=none,numbers=none,mathescape=true,escapechar=']
requires Licensor_balance >= 0 && Licensee_balance > 0 &&
         Licensor_token > 0    && Licensee_token == 0 ;
ensures  (Licensor_balance == \old(Licensor_balance)+cost && Licensor_token == 0 
          && Licensee_balance == \old(Licensee_balance)-cost 
          && Licensee_token == \old(Licensor_token)
         ) || ( 
          Licensor_balance == \old(Licensor_balance) &&
          Licensor_token == \old(Licensor_token) && 
          Licensee_balance == \old(Licensee_balance) &&
          Licensee_token == \old(Licensee_token)
         ) ;
\end{lstlisting}
}

The contract expresses that, upon termination, either the licensing
transaction succeeded (the \stipula{Licensee} purchased the license
and the \stipula{Licensor} obtained the money) or the
\stipula{Licensee} is withdrawn and no \stipula{token} was delivered.

\begin{theorem}[Soundness]
  \label{thm.soundnessacyclic}
  Let $\C$ be an acyclic {\Stipula} contract
  % Let $\contract {\comma} {\tt 0}$ be the initial configuration of a
  % {\Stipula} acyclic contract {\C}
  and let
  \[
    \contract {\comma} {\tt 0} \lred{(\vect A,\, \vect{A_i}{:} \vect{v_i}{\,}^{i \in 1..\ell})}
    \lred{}^* \contract_1 {\comma} \Time_1 \lred{\mu_1} \lred{}^* 
    \contract_2 {\comma} \Time_2 \lred{\mu_2} \lred{}^* \cdots 
    % \contract_{n} {\comma} \Time_n
    \lred{\mu_{n+1}} \lred{}^* \contract_{n+1} {\comma} \Time_{n+1}
  \]
  where $\Q_i$ are the states of $\contract_i$, with $i \in 1..n+1$,
  $\Q_{n+1}$ is the final state of the underlying automaton, and
  $\mu_i$ are either function invocations
  $A_i{:}\f_i(\vect{v_i})[\vect{v_i'}]$ or events $\ev_i$. Then the
  above computation may be traced in the code of
  $\genscenario{\C}{\cdot}$ as follows. Consider
  $\contract_i {\comma} \Time_i \lred{\mu_i} \lred{}^* \contract_{i+1}
  {\comma} \Time_{i+1}$; if, in the code of $\genscenario{\C}{\Q_i}$,
  \begin{enumerate}
  \item \stipula{now} = $\Time_i$ if $\mu_i$ is an event;
    \stipula{now} $\leq \Time_i$ if $\mu_i$ is a function;
  \item if the multiset of pending events in $\contract_i$ is
    \[
      {\Time}_{\jmath_1} \,\event_{\ev_{\jmath_1}}\,
      \Qwithat_{\jmath_1}\, \{ \, S_{\jmath_1} \,\} \,\tostate\,
      \Qwithat_{\jmath_1}' ~|~ \cdots ~|~ {\Time}_{\jmath_\kappa}
      \event_{\ev_{\jmath_k}}\, \Qwithat_{\jmath_k}\, \{ \,
      S_{\jmath_k} \,\} \,\tostate\, \Qwithat_{\jmath_k}'
    \]
    and the dispatch table
    ${\tt DT}[\ev_{\jmath_i}] = {\Time}_{\jmath_i}$, with
    $i \in 1.. k$.
\end{enumerate}
Then, whenever $\mu_i = A_i{:}\f_i(\vect{v_i})[\vect{v_i'}]$ it is
possible to execute the method $\f_i(\vect{v_i},\vect{v_i'})$ and if
$\mu_i = \ev_{\jmath_i}$, it is possible to execute method
${\tt event}\_{\jmath_i}(~)$ and obtain a state where (1.) and (2.)
hold for $\contract_{i+1} \comma \Time_{i+1}$.
\end{theorem}

Therefore, every computation of an acyclic {\Stipula} contract is
faithfully represented within the {\Java} code generated by our
translation
%
% for each execution permitted by the {\Stipula} semantics of the
% contract, there exists a corresponding execution path in the
% translated {\Java}/{\JML} program.
(the proof is omitted, it is technical, but not difficult).
Consequently, when the program is analysed with {\KeY}, each legally
admissible behaviour of the original contract is systematically
explored. This guarantees that the verification process is sound with
respect to the contract's operational semantics in the acyclic
setting, as no admissible execution is omitted.  We further conjecture
that the translation is complete in the acyclic case, in the sense
that every behaviour of the generated {\Java} program corresponds to a
computation that is admissible in the original {\Stipula}
contract. Establishing such a result formally, however, would require
a formalisation of {\Java} semantics, which lies beyond the
scope of the present paper. 

%%% Local Variables:
%%% mode: latex
%%% TeX-master: "newmain"
%%% End:


\section{Cyclic Behaviour and its Translation}
\label{sec:cyclic-behavior}
%!TEX root = newmain.tex

  The technique of Sections~\ref{sec:translation}
  and~\ref{sec:dispatchtable} extends to cyclic contracts with
  separate cycles such as \stipula{Deposit}. Here we examine the
  additional challenges posed by cycles and clarify the resulting loss
  of precision.

  In \stipula{Deposit}, the \stipula{Farm} and the \stipula{Client}
  may repeatedly invoke \stipula{send} and \stipula{buy}, forming a
  cycle in the contract underlying automaton
  (Figure~\ref{fig:deposit-fsm}). The cycle is exited through a
  dedicated event executed at most once. This reflects a common
  contractual pattern: (\emph{i}) cycles have a unique entry/exit
  state and contain only function invocations; (\emph{ii}) cycles are
  not nested; (\emph{iii}) events occur outside cycles and typically
  enforce their termination. As a consequence, each event is executed
  at most once, even in the presence of cyclic behaviour.
  
  This observation allows us to reuse the dispatch table mechanism
  developed for acyclic contracts, since the set of events remains
  finite.  However, events generated inside cycles require special
  care: although they share the same index, different iterations may
  assign them different time values.  To handle this, we assign such
  events an \emph{abstract} clock value, encoded as \jml{-2}.  Rather
  than reasoning about their precise scheduling time, we only record
  that they have been generated. We illustrate the resulting loss
  of precision with an example. Consider the cyclic \Stipula\ function
  declaration:
  
  {\small
\begin{lstlisting}[numbers=none,frame=none,mathescape=true,language=stipula,escapechar=§]
@Q f()[] {
   now + 1 $\event$ @Q {} $\tostate$ @Q1
   now + 2 $\event$ @Q {} $\tostate$ @Q2
} $\tostate$ @Q
\end{lstlisting}
  }

  According to the \Stipula\ semantics, the second event is never
  triggered, because the first one always precedes it, so state
  \stipula{@Q2} is unreachable. But since the events occur in a cycle,
  both are given the abstract clock value \jml{-2}, enabling an
  infeasible execution in the \JML/\Java\ model that triggers the
  second event and reaches \stipula{@Q2}.

  For modeling cycles in \JML\ we refine $\updateDT{\Q}{W}$ and
  distinguish whether the generating function of an event lies inside
  or outside a cycle. If it belongs to a cycle, the events in $W$ are
  inserted into the dispatch table with the abstract time value
  \jml{-2}, instead of a concrete timestamp.  Accordingly, the
  predicate \ensureDT$_{\f}$ is adapted to enforce that \emph{exactly}
  the events generated by $\f$ are assigned the value \jml{-2}, while
  all other entries of \jml{DT} remain unchanged. The structure of the
  specification remains the same as in the acyclic case, differing
  only in the assigned time value. The remaining definitions in
  Section~\ref{sec:dispatchtable} are omitted, because they are
  straightforward.

  The crucial adaptation concerns the extension of
  $\genscenario{\C}{\Q}$ (Listing~\ref{tab:fungenerate}) to states
  within cycles. For such states, we use the notation
  $\genscenarioplus{\C}{\Q,{\tt count}}{\Q'}$, where $\Q$ denotes the
  entry/exit state of the cycle and \jml{count} is the iteration
  counter introduced by the translation (see
  Listing~\ref{tab:fungeneratecycles}).  Each cycle is translated into
  a dedicated {\Java} \jml{while} loop whose body encodes all possible
  traversals of the cycle. A counter tracks the number of iterations,
  bounded by a cycle-specific constant \stipula{MAX_ITE}$_\Q$. This
  bound becomes a parameter of the \stipula{behaviour} method in
  Listing~\ref{tab:translate} with precondition %
  \jml{requires MAX_ITE$_\Q\geq\,$0} and provides a natural
  abstraction: although {\Stipula} allows unbounded cyclic behaviour
  (\emph{e.g.}, in \stipula{Deposit}), only finitely many interactions
  occur in practice.  Finally, the time-progress variable
  \stipula{u_Q} used in the acyclic case is omitted for cyclic states,
  since events generated within cycles carry an undetermined time
  value and do not require explicit clock advancement.

{\small
\begin{lstlisting}[frame=none,numbers=none,language={[JML]Java},mathescape=true,caption={The function $\genscenarioplus{\C}{\Q,{\tt count}}{\Q'}$.},label={tab:fungeneratecycles},float,escapechar=�]
�{\textrm{There are two cases:\\
  \hspace*{-2em}(1) $\Q$ is the entry/exit state of a cycle; ${\f}_1, \ldots, {\f}_n$ and\ 
  $\mathtt{ev_{\jmath\,1}}, \ldots, {\tt ev_{\jmath\,m}}$ are functions\\
  \hspace*{-2em} with initial state $\Q$ and $\texttt{final}$ states\ 
  $\Q_1, \ldots, \Q_n$ and $\Q_1', \ldots, \Q_m'$, respectively.}�

$\genscenarioplus{\C}{}{\Q}  \; \eqdef$
   int entry = minTimeEntry({ev$_{\jmath 1}$, $\cdots$ , ev$_{\jmath m}$}) ; 
   if ((entry == -1) || (DT[entry] > now && w_Q > 0)) {
     int count = 0 ;
     /*@ loop_invariant 0 <= count <= MAX_ITE$_\Q$ ; // bounds
       @ loop_invariant ... ; // assets, fields, and DT invariants     
       @ decreases MAX_ITE$_\Q$ - count ; // termination expression
       @*/
     while (count < MAX_ITE$_\Q$) {
       if (w_Q == 1) { $\f_1(\vect{\tt x_1}, \vect{\tt k_1})$ ; if (Q != Q$_1$) $\genscenarioplus{\C}{\Q,{\tt count}}{\Q_1}$ }
       else if (w_Q == 2) { $\f_2(\vect{\tt x_2}, \vect{\tt k_2})$ ; if (Q != Q$_2$) $\genscenarioplus{\C}{\Q,{\tt count}}{\Q_2}$ }
       ...
       else { $\f_n(\vect{{\tt x}_n}, \vect{{\tt k}_n})$ ;  if (Q != Q$_n$) $\genscenarioplus{\C}{\Q,{\tt count}}{\Q_n}$ }  // max_Q == n
       count = count + 1 ;
     } 
     entry = minTimeEntry({ev$_{\jmath 1}$,$\cdots$,ev$_{\jmath m}$}) ; // we are back in state Q
     if (DT[entry] > now || DT[entry] == -2) { 
       if (DT[entry] > now) now = DT[entry] ;
       if (entry == ev$_{\jmath 1}$) { event_${\jmath}$1() ; $\genscenarioplus{\C}{}{\Q_1'}$ }
       else if (entry == ev$_{\jmath 2}$) { event_${\jmath}$2(); $\genscenarioplus{\C}{}{\Q_2'}$ }
       ...
       else { event_${\jmath}$m() ; $\genscenarioplus{\C}{}{\Q_m'}$ }      //  entry == ev$_{\jmath m}$
     } 
   } else { if (DT[entry] > now) now = DT[entry] ;
     if (entry == ev$_{\jmath 1}$) { event_${\jmath}$1() ; $\genscenarioplus{\C}{}{\Q_1'}$ }
     else if (entry == ev$_{\jmath 2}$) { event_${\jmath}$2(); $\genscenarioplus{\C}{}{\Q_2'}$ }
     ...
     else { event_${\jmath}$m() ; $\genscenarioplus{\C}{}{\Q_m'}$ }      //  entry == ev$_{\jmath m}$
   }    
     
 �\textrm{(2) ${\Q}'$ is a state of a cycle that is \emph{not} the entry/exit state. Let ${\f}_1, \ldots, {\f}_\ell$ be the functions whose final states $\Q_1, \ldots, \Q_\ell$ are in the cycle.}�

$\genscenarioplus{\C}{\Q,{\tt count}}{\Q'} \eqdef$ if (w_Q$'$ == 1) { $\f_1(\vect{\tt x_1},\vect{\tt k_1})$ ; if (Q != Q$_1$) $\genscenarioplus{\C}{\Q,{\tt count}}{\Q_1}$ }
                else if (w_Q$'$ == 2) { $\f_2(\vect{\tt x_2}, \vect{\tt k_2})$ ; if (Q != Q$_2$) $\genscenarioplus{\C}{\Q,{\tt count}}{\Q_2}$ }
                ...
                else { $\f_\ell(\vect{\tt x_\ell}, \vect{\tt k_\ell})$ ; if (Q != Q$_\ell$) $\genscenarioplus{\C}{\Q,{\tt count}}{\Q_\ell}$ } 
\end{lstlisting}
}

From a deductive verification perspective,
% using counters to define loop iterations has an important advantage:
% it simplifies the generation of loop invariants in {\KeY}.  In
% practice, the invariant of \jml{while} loop can relate the current
% state of the contract, the dispatch table, and the iteration
% counter, thus reducing the burden of invariant discovery and
% enabling {\KeY} to reason effectively about cyclic behaviors within
% a bounded yet sound approximation.  Of course,
finding suitable loop invariants is a very difficult task in formal
program verification and usually not amenable to automation. For this
reason our translation technique leaves the supply of post
condition-specific invariants to the verification engineer.
%
For example, in the case of \stipula{Deposit}, using the following
invariant ($v_s$ is the actual value of the parameters {\tt h} and
{\tt k} in {\tt begin} and {\tt send}; $v_b$ is the value of {\tt w}
in {\tt buy}) it is possible to derive in {\KeY} the property that all
the flour acquired by the \stipula{Client} has been paid to the
\stipula{Farm} at agreed cost.

{\small
\begin{lstlisting}[frame=none,numbers=none,language={[JML]Java},mathescape=true]
loop_invariant (Client_flour == \old(Client_flour) && flour == \old(flour)+$v_{s}$)$\,$||
               (Client_flour == \old(Client_flour)+$v_b$/cost_flour &&
                Farm == \old(Farm)+$v_b$) ;
\end{lstlisting}   }   

  Another limitation of our current encoding is that function
  parameters are treated as symbolic constants throughout a loop,
  meaning all iterations are analysed under a single symbolic
  instantiation (e.g., $v_s$ and $v_b$). While {\Stipula} allows
  parameters to vary between iterations, verifying this
  non-determinism would require loop invariants over sequences of
  iteration-dependent values, which tools like {\KeY} cannot currently
  handle automatically. Supporting such cases is theoretically
  possible with user-supplied invariants and interactive proofs, but
  will further reduce the automation of our approach.

For the cyclic case, we can prove a soundness theorem for the contracts analysed here, 
assuming constraints on repeated function calls with identical arguments and on 
the number of iterations.

\begin{theorem}[Soundness]
  \label{thm.soundnessacyclic}
  Let $\C$ be a {\Stipula} contract
  % $\contract {\comma} 0$ be the initial configuration of a
  % {\Stipula} contract {\C}
  with separate cycles and events that are outside cycles. Let
  \[
    \contract {\comma} {\tt 0} \lred{(\vect A,\, \vect{A_i}{:} \vect{u_i}{\,}^{i \in 1..\ell})}
    \lred{}^* \contract_1 {\comma} \Time_1 \lred{\mu_1} \lred{}^* 
    \contract_2 {\comma} \Time_2 \lred{\mu_2} \lred{}^* \cdots 
    % \contract_{n} {\comma} \Time_n
    \lred{\mu_{n+1}} \lred{}^* \contract_{n+1} {\comma} \Time_{n+1}
  \]
  where $\Q_i$ are the states of $\contract_i$, with $i \in 1..n+1$,
  and $\mu_i$ are either function invocations
  $A_i{:}\f_i(\vect{v_i})[\vect{v_i'}]$ or events
  $\ev_i$. Additionally, if $\mu_i$ and $\mu_j$ are invocations of the
  same function then $\mu_i = \mu_j$ (the invocations have the same
  actual values).
  % 
  Then, %with suitable values of the constants \stipula{MAX_ITE}$_\Q$
  the above computation may be traced in the code of
  $\genscenario{\C}{\cdot}$ as follows. Consider
  $\contract_i {\comma} \Time_i \lred{\mu_i} \lred{}^* \contract_{i+1}
  {\comma} \Time_{i+1}$ and assume that, in the code of
  $\genscenario{\C}{\Q_i}$ or of
  $\genscenarioplus{\C}{\Q,{\tt count}}{\Q_i}$ (when $\Q_i$ is part of
  a cycle),
  \begin{enumerate}
  \item \stipula{now} = $\Time_i$ if $\mu_i$ is an event not generated
    inside a cycle; \stipula{now} $\leq \Time_i$, otherwise;
  \item if the multiset of pending events in $\contract_i$ is
    \[
      {\Time}_{\jmath_1} \,\event_{\ev_{\jmath_1}}\,
      \Qwithat_{\jmath_1}\, \{ \, S_{\jmath_1} \,\} \,\tostate\,
      \Qwithat_{\jmath_1}' ~|~ \cdots ~|~ {\Time}_{\jmath_\kappa}
      \event_{\ev_{\jmath_k}}\, \Qwithat_{\jmath_k}\, \{ \,
      S_{\jmath_k} \,\} \,\tostate\, \Qwithat_{\jmath_k}'
    \]
    and the dispatch table is such that
    ${\tt DT}[\ev_{\jmath_i}] \in \{ {\Time}_{\jmath_i}, \mbox{\tt
      -2}\}$, with $i \in 1.. k$.
  \end{enumerate}
  Then, if $\mu_i = A_i{:}\f_i(\vect{v_i})[\vect{v_i'}]$ it is
  possible to execute the method $\f_i(\vect{v_i},\vect{v_i'})$ and if
  $\mu_i = \ev_{\jmath_i}$, it is possible to execute method
  ${\tt event}\_{\jmath_i}(~)$ and obtain a state where (1.) and (2.)
  hold for $\contract_{i+1} \comma \Time_{i+1}$.
\end{theorem}

This result crucially depends on the values of \stipula{MAX_ITE}$_\Q$.
If reliable upper bounds on the number of possible iterations are
known, the analysis carried out with {\KeY} is guaranteed to explore
all the legally admissible behaviour of the contract. However, in
contrast to the acyclic setting, completeness can no longer be
claimed.  As discussed at the beginning of the section, the introduction of over-approximations, specifically,
assigning the abstract time value {\tt -2} to events generated inside
cycles, may lead the translated {\Java} program to exhibit behaviour
not admissible by the semantics of {\Stipula}. As a consequence, the
verification performed by {\KeY} may encounter spurious executions,
potentially preventing the proof of properties that do hold in the
original contract.

%%% Local Variables:
%%% mode: latex
%%% TeX-master: "newmain"
%%% End:


%\section{Implementation and Evaluation}
%\label{sect:casestudies}
%
\section{Related Work}
\label{sec:relatedworks}

Legal modeling frameworks like Ergo~\cite{ergo},
OpenLaw~\cite{OpenLaw}, Lexon~\cite{Lexon}, and Accord~\cite{Accord}
embed contracts into broader systems, but they lack a precise formal
semantics. While Catala~\cite{Catala} formalizes legislative logic and
has a runtime environment, it does not address contract
verification. In contrast, {\Stipula} offers an operational model with
explicit permissions, assets, and timed clauses and has a
comprehensive suite of tools, including a visual code editor for
intuitive contract authoring, an interpreter for automatic execution,
and a number of analyzers for verifying several legally relevant
properties (unreachability of clauses, absence of frozen assets,
etc.)~\cite{Laneve2026}.

Closer to our work, prior efforts explored translation-based
verification: OCL-to-{\Java} with {\JML}~\cite{hamie2002ocl}, and
Circus-to-{\Java} for formal reasoning~\cite{DBLP:conf/fm/FreitasC06}.
Our approach applies this paradigm to legal contracts, preserving
their normative and temporal semantics.

More recent work models blockchain smart contracts (by hand) in a
modeling language designed for verification which is then
automatically translated to EVM bytecode \cite{Solidity}, for example,
Dafny-to-EVM \cite{Dafny24} or Why3-to-EVM \cite{Why319}. This is the
\emph{reverse} direction of our approach. While it makes verification
easier, \emph{validation} becomes harder.

While the current {\Stipula} runtime is realized by compilation to
{\Java}~\cite{stipulaprototype}, the language is
implementation-agnostic and could be compiled to blockchain smart
contracts like {\tt Solidity}~\cite{Solidity} or {\tt
  Obsidian}~\cite{Obsidian}.  This path is promising, as our
verification approach using {\KeY}~\cite{DBLP:series/lncs/10001} can
build on established work that already applies the system to
blockchain platforms like Hyperledger
Fabric~\cite{BeckertHerdaKirsten2018_1000092715} and {\tt
  Solidity}~\cite{DBLP:conf/isola/AhrendtB20} by embedding them into
\Java.
% Using these advances would enable rigorous verification of
% {\Stipula} contracts on decentralized infrastructures.


\section{Conclusion}
\label{sec:conclusion}

  The formal modeling and verification of digital contracts is
  gaining momentum at the intersection of legal informatics,
  programming languages, and formal methods.  This work contributes by
  connecting the contract language {\Stipula} with a deductive
  verification system (\KeY) via a translation to {\JML}-annotated
  {\Java}.
 
  We focus on two recurring interaction patterns in legal contracts:
  acyclic behaviour and non-nested cyclic behaviour consisting only of
  function calls. We also impose other technical restrictions which
  might be lifted in the future.  Maximising automation has been our
  guiding principle. In the acyclic case, verification is fully
  automatic once the contract is specified. For cyclic behaviour, where
  stronger reasoning is needed, we aim to minimise additional
  annotations and proof guidance to preserve usability and
  scalability. From a theoretical perspective, the most challenging
  problem is the precise handling of clock values in events
  originating from cycles. It would require a generalisation of
  dispatch tables that precisely describes event generation for an
  \emph{unbounded} number of iterations.
  %
  A systematic, experimentally driven investigation of
  the trade-offs between coverage and precision constitutes an
  important direction for future work.  Table~\ref{tab:restrictions}
  summarises the current restrictions and to which extent they could
  be dropped.

  \begin{table}[t]
\begin{center}
{\small
  \begin{tabular}{p{.43\textwidth}@{\quad\;}p{.52\textwidth}@{}}
    \hline  
    {\bf Restriction} & {\bf Motivation / How to Lift}\\
    \hline
    Fields, assets and variables are of type integer, see~Section~\ref{subsect:stipula}
    &
    Increase degree of automation; can be lifted by dedicated automated reasoners
    \\[1ex]
%    In time guards \stipula{now+t}, time values \stipula{t} are constants
%    &
%    Symbolic values in dispatch tables increase branching degree of symbolic execution: lifting subject to experiments
%    \\[1ex]
    Cycles in contracts are separate and event-free, see~Sections~\ref{subsect:stipula} and~\ref{sec:cyclic-behavior}
    &
    General cycles are rare in legal contracts and significantly complicate the translation, because they cannot be modeled by loops
    \\[1ex]
    Events generated in cycles get abstract clock value, see Section~~\ref{sec:cyclic-behavior}
    &
    Precise treatment requires dynamic and unbounded extension of dispatch table
    \\[1ex]  
    Cycle iterations are analysed under a single symbolic instantiation of the formal parameters, see
    Section~\ref{sec:cyclic-behavior}
    &
    Requires more complex loop invariants that might be hard to find\\
    \hline
  \end{tabular}
}
\end{center}
\caption{\label{tab:restrictions} Restrictions to {\Stipula} and motivations} 
\vspace{-.5cm}
\end{table}

%%% Local Variables:
%%% mode: latex
%%% TeX-master: "newmain"
%%% End:

  
  % Some technical restrictions could be relaxed: for
  % example, allowing events in cycles when states do not mix functions
  % and events, but others are more fundamental. In particular,
  % supporting nested cycles or time expressions depending on parameters
  % or contract fields would require substantial changes to the
  % translation, complicating the execution-path encoding and dispatch
  % management, and potentially reducing the automation and strength of
  % the verification in {\KeY}.

  In general, a systematic empirical evaluation, especially for large
  and structurally complex cyclic contracts remains to be
  done. Preliminary results are promising but concern an earlier
  version of the translation without dispatch tables~\cite{iFM25}.  A
  comprehensive assessment of the current framework remains essential.


\appendix
\vspace{-.3cm}
\section{Addendum: The Function {\sc Stipula2Java}}
\label{sec:stipula2java}
%!TEX root = newmain.tex

The definition of $\StipulaToJava{\A}{S}$ follows.

\vspace{-.1cm}

\begin{itemize}
\item[--] $\StipulaToJava{\A}{ \zero} \eqdef \varepsilon$, where
  $\varepsilon$ is the empty string.
%\medskip
\item[--]
  $\StipulaToJava{\A}{E \, \send \, {\x} \; \; S} \eqdef \; \, \x \,
  \mbox{\tt =} \, E \ttsemi \StipulaToJava{\A}{ S}$.
%\medskip
\item[--]
  $\StipulaToJava{\A}{E \, \send \, {\B} \; \; S} \eqdef
  \StipulaToJava{\A}{ S}$. The value of $E$ is irrelevant in our
  examples. If the value of $E$ is relevant, we declare an \emph{ad
    hoc} field, say \stipula{B_transmitted}, for storing the value.
%\medskip
\item[--] $\StipulaToJava{\A}{E \, \lolli \, \h, \h' \; \; S}$. There
  are two cases: (\emph{i}) both $\h$ and $\h'$ are assets of the
  contract or 
  
  (\emph{ii}) $\h$ is a formal parameter and $\h'$ is an
  asset of the contract. The code for case (\emph{i}) and (\emph{ii})
  is ({\tt tmp} is a fresh variable):
  \begin{center}
    {\small
      \begin{tabular}{c@{\qquad\quad}c}
        $
        \begin{array}{l}
          {\tt int \; tmp} \, \mbox{\tt =} \, E \ttsemi 
          \\
          \mbox{\jml{//@ $\;$ assert h-tmp >= 0}} \ttsemi 
          \\
          \h \, \mbox{\tt =} \, \h \mbox{\tt -} {\tt tmp} 
          \ttsemi \h' \, \mbox{\tt =} \, \h' \mbox{\tt +} {\tt tmp} \ttsemi 
          \\
          \StipulaToJava{\A}{ S}
        \end{array}
        $
        &
          $
          \begin{array}{l}
            {\tt int \; tmp} \, \mbox{\tt =} \, E \ttsemi 
            \\
            \mbox{\jml{//@ $\;$ assert h-tmp >= 0  \&\& \A${\unds}$\h'-tmp >= 0}} \ttsemi 
            \\
            \h \, \mbox{\tt =} \, \h \mbox{\tt -} {\tt tmp} 
            \ttsemi \A{\unds}\h' \, \mbox{\tt =} \, \A{\unds}\h' \mbox{\tt -} {\tt tmp} 
            \ttsemi
            \\
            \h' \, \mbox{\tt =} \, \h' \mbox{\tt +} {\tt tmp} \ttsemi 
            \\
            \StipulaToJava{\A}{ S}
          \end{array}
          $
        \\[.2cm]
        Case \emph{(i)} & Case \emph{(ii)}
      \end{tabular}}
  \end{center}

\noindent
% In case {\tt tmp} is a fresh identifier storing the value of the
% expression $E$.
According to the \JML\ \jml{assert} semantics, verification can only
succeed provided that \jml{h-tmp >= 0}, which is exactly what the
{\Stipula} semantics requires. In addition, in case (\emph{ii}), the
formal parameter {\h} represents (part of) the asset {\A}{\tt
  \_}{\h}$'$ that the caller {\A} transfers to the contract upon
invocation. Consequently, the decrease in {\tt h} must be reflected in
the caller's asset rather than in the contract state. This semantic
correspondence motivates the assignment {\tt A{\unds}h$'$ =
  A{\unds}h$'$-tmp;} which correctly accounts for the transfer of
the asset from the caller to the contract.

\medskip

\item[--] $\StipulaToJava{\A}{E \, \lolli \, \h, \B \; \; S}$.  Again,
  there are two cases. Either (\emph{i}) {\h} is an asset or
  (\emph{ii}) it is a formal parameter. The code of cases (\emph{i})
  and (\emph{ii}) is:
  %
  \begin{center}
    {\small
      \begin{tabular}{c@{\qquad\quad}c}
        $\begin{array}{l}
           {\tt int \; tmp} \, \mbox{\tt =} \, E \ttsemi 
           \\
           \mbox{\jml{//@ $\;$ assert h-tmp >= 0}}  \ttsemi 
           \\
           \h \, \mbox{\tt =} \, \h \mbox{\tt -} {\tt tmp} \ttsemi 
           \\
           \B{\unds}\h \, \mbox{\tt =} \, \B{\unds}\h \mbox{\tt +} {\tt tmp} \ttsemi 
           \\
           \StipulaToJava{\A}{ S}
         \end{array}$
        &
          $\begin{array}{l}
             {\tt int \; tmp} \, \mbox{\tt =} \, E \ttsemi 
             \\
             \mbox{\jml{//@ $\;$ assert h-tmp >= 0 \&\& \A-tmp >= 0}}  \ttsemi 
             \\
             \h \, \mbox{\tt =} \, \h \mbox{\tt -} {\tt tmp} 
             \ttsemi \A \, \mbox{\tt =} \, \A \mbox{\tt -} {\tt tmp} \ttsemi 
             \\
             \B \, \mbox{\tt =} \, \B \mbox{\tt +} {\tt tmp} \ttsemi 
             \\
             \StipulaToJava{\A}{ S}
           \end{array}$
        \\[.2cm]
        Case \emph{(i)} & Case \emph{(ii)}
      \end{tabular}}
  \end{center}

\noindent
In case (\emph{ii}), we decrease and increase the assets {\A} and
{\B}, respectively, where {\tt tmp} is a fresh variable. In case
(\emph{i}), we assume that $\h$ is an asset, not a formal
parameter. This allows us to determine what asset of {\tt B} must be
increased by the move (the general case is dealt with \emph{ad hoc}
annotations).

\medskip

\item $\StipulaToJava{\A}{\ifte{E}{S'}{S''} \; S} \eqdef \; \,$
    {\small
      $\begin{array}{l} {\tt if} \, \mbox{\tt (} \, E \, \mbox{\tt )
         \{} \; \StipulaToJava{\A}{S'} \; \mbox{\tt \}}
         \\
         {\tt else} \; \mbox{\tt \{} \; \StipulaToJava{\A}{ S''} \;
         \mbox{\tt \}}
         \\
         \StipulaToJava{\A}{ S}
       \end{array}$
     }
     
\medskip

\noindent
The semantics of the conditional is obvious: we simply use the
corresponding {\Java} statement.
\end{itemize} 

Regarding the shortcuts $\h \lolli \h'$ and $\h \lolli \A$, we
generate optimized {\Java} code, to avoid the declaration of the
auxiliary variable {\tt tmp}:
\[
  \StipulaToJava{\A}{\h \lolli \h'} \eqdef \; {\small \jml{h' = h'+h ; h = 0 ;}}\ \mbox{\textcolor{commentColor}{{\small {\tt // \h \; and \h' assets}}}}
\]
\[
  \StipulaToJava{\A}{\h \lolli \h'} \eqdef \;
  {\small
    \begin{array}{ll} 
      \mbox{\jml{//@ $\;$ assert \A${\unds}$h'-h >= 0}} \ttsemi & 
                                                                   \mbox{\textcolor{commentColor}{{\small {\tt // \h\ formal param., \h' asset}}}}
      \\
      \jml{h' = h'+h ;} & 
      \\
      \jml{A_h' = A_h'-h ;} & 
      \\
      \jml{h = 0 ;} &
    \end{array}
  }
\]
%and 
\[
\begin{array}{l@{\quad}l}
{\rm and} & \StipulaToJava{\A}{\h \lolli \B} \eqdef \; \jml{B_h = B_h+h ; h = 0 ;}\ \mbox{\textcolor{commentColor}{{\small {\tt // \h \; asset}}}} 
\\[.2cm]
& \StipulaToJava{\A}{\h \lolli \B} \eqdef \; 
\begin{array}{l}
\mbox{\jml{//@ $\;$ assert \A-h >= 0}} \ttsemi\ \mbox{\textcolor{commentColor}{{\small {\tt // \h \; formal parameter}}}} 
\\
\jml{B = b+h ; h = 0 ;}
\end{array}
\end{array}
\]

%%% Local Variables:
%%% mode: latex
%%% TeX-master: "newmain"
%%% End:


\bibliographystyle{splncs04}
\bibliography{references}


\end{document}

%%% Local Variables:
%%% mode: latex
%%% TeX-master: t
%%% End:
